%%
%% This is file `t-markdown.tex',
%% generated with the docstrip utility.
%%
%% The original source files were:
%%
%% markdown.dtx  (with options: `context')
%% 
%% Copyright (C) 2016-2022 Vít Novotný
%% 
%% This work may be distributed and/or modified under the
%% conditions of the LaTeX Project Public License, either version 1.3c
%% of this license or (at your option) any later version.
%% The latest version of this license is in
%% 
%%    http://www.latex-project.org/lppl.txt
%% 
%% This work has the LPPL maintenance status `maintained'.
%% The Current Maintainer of this work is Vít Novotný.
%% 
%% Send bug reports, requests for additions and questions
%% either to the GitHub issue tracker at
%% 
%%   https://github.com/Witiko/markdown/issues
%% 
%% or to the e-mail address <witiko@mail.muni.cz>.
%% 
%% MODIFICATION ADVICE:
%% 
%% If you want to customize this file, it is best to make a copy of
%% the source file(s) from which it was produced.  Use a different
%% name for your copy(ies) and modify the copy(ies); this will ensure
%% that your modifications do not get overwritten when you install a
%% new release of the standard system.  You should also ensure that
%% your modified source file does not generate any modified file with
%% the same name as a standard file.
%% 
%% You will also need to produce your own, suitably named, .ins file to
%% control the generation of files from your source file; this file
%% should contain your own preambles for the files it generates, not
%% those in the standard .ins files.
%% 
%% The names of the source files used are shown above.
%% 
\writestatus{loading}{ConTeXt User Module / markdown}%
\startmodule[markdown]
\unprotect
\def\markdownInfo#1{\writestatus{markdown}{#1.}}%
\def\markdownWarning#1{\writestatus{markdown\space warn}{#1.}}%
\def\dospecials{\do\ \do\\\do\{\do\}\do\$\do\&%
  \do\#\do\^\do\_\do\%\do\~}%
\input markdown/markdown
\let\startmarkdown\relax
\let\stopmarkdown\relax
\let\inputmarkdown\relax
\ExplSyntaxOn
\cs_new:Nn
  \__markdown_setup:n
  {
    \keys_set:nn
      { markdown/context-options }
      { #1 }
  }
\long\def\setupmarkdown[#1]
  {
    \__markdown_setup:n
      { #1 }
  }
\ExplSyntaxOff
\ExplSyntaxOn
\cs_new:Nn \__markdown_context_define_option_commands_and_keyvals:
  {
    \seq_map_inline:Nn
      \g__markdown_option_layers_seq
      {
        \seq_map_inline:cn
          { g__markdown_ ##1 _options_seq }
          {
              \__markdown_context_define_option_keyval:nn
                { ##1 }
                { ####1 }
          }
      }
  }
\cs_new:Nn \__markdown_context_define_option_keyval:nn
  {
    \prop_get:cnN
      { g__markdown_ #1 _option_types_prop }
      { #2 }
      \l_tmpa_tl
    \keys_define:nn
      { markdown/context-options }
      {
        #2 .code:n = {
          \tl_set:Nx
            \l_tmpa_tl
            {
              \str_case:nnF
                { ##1 }
                {
                  { yes } { true }
                  { no } { false }
                }
                { ##1 }
            }
          \__markdown_set_option_value:nV
            { #2 }
            \l_tmpa_tl
        },
      }
    \str_if_eq:VVT
      \l_tmpa_tl
      \c__markdown_option_type_boolean_tl
      {
        \keys_define:nn
          { markdown/context-options }
          {
            #2 .default:n = { true },
          }
      }
  }
\cs_generate_variant:Nn
  \__markdown_set_option_value:nn
  { nV }
\__markdown_context_define_option_commands_and_keyvals:
\ExplSyntaxOff
\def\markdownMakeOther{%
  \count0=128\relax
  \loop
    \catcode\count0=11\relax
    \advance\count0 by 1\relax
  \ifnum\count0<256\repeat
  \catcode`|=12}%
\long\def\inputmarkdown{%
  \dosingleempty
  \doinputmarkdown}%
\long\def\doinputmarkdown[#1]#2{%
  \begingroup
    \iffirstargument
      \setupmarkdown{#1}%
    \fi
    \markdownInput{#2}%
  \endgroup}%
\ifx\startluacode\undefined  % MkII
  \begingroup
    \catcode`\|=0%
    \catcode`\\=12%
    |gdef|startmarkdown{%
      |markdownReadAndConvert{\stopmarkdown}%
                             {|stopmarkdown}}%
    |gdef|stopmarkdown{%
      |markdownEnd}%
  |endgroup
\else  % MkIV
  \startluacode
    document.markdown_buffering = false
    local function preserve_trailing_spaces(line)
      if document.markdown_buffering then
        line = line:gsub("[ \t][ \t]$", "\t\t")
      end
      return line
    end
    resolvers.installinputlinehandler(preserve_trailing_spaces)
  \stopluacode
  \begingroup
    \catcode`\|=0%
    \catcode`\\=12%
    |gdef|startmarkdown{%
      |ctxlua{document.markdown_buffering = true}%
      |markdownReadAndConvert{\stopmarkdown}%
                             {|stopmarkdown}}%
    |gdef|stopmarkdown{%
      |ctxlua{document.markdown_buffering = false}%
      |markdownEnd}%
  |endgroup
\fi
\def\markdownRendererLineBreakPrototype{\blank}%
\def\markdownRendererLeftBracePrototype{\textbraceleft}%
\def\markdownRendererRightBracePrototype{\textbraceright}%
\def\markdownRendererDollarSignPrototype{\textdollar}%
\def\markdownRendererPercentSignPrototype{\percent}%
\def\markdownRendererUnderscorePrototype{\textunderscore}%
\def\markdownRendererCircumflexPrototype{\textcircumflex}%
\def\markdownRendererBackslashPrototype{\textbackslash}%
\def\markdownRendererTildePrototype{\textasciitilde}%
\def\markdownRendererPipePrototype{\char`|}%
\def\markdownRendererLinkPrototype#1#2#3#4{%
  \useURL[#1][#3][][#4]#1\footnote[#1]{\ifx\empty#4\empty\else#4:
  \fi\tt<\hyphenatedurl{#3}>}}%
\usemodule[database]
\defineseparatedlist
  [MarkdownConTeXtCSV]
  [separator={,},
   before=\bTABLE,after=\eTABLE,
   first=\bTR,last=\eTR,
   left=\bTD,right=\eTD]
\def\markdownConTeXtCSV{csv}
\def\markdownRendererContentBlockPrototype#1#2#3#4{%
  \def\markdownConTeXtCSV@arg{#1}%
  \ifx\markdownConTeXtCSV@arg\markdownConTeXtCSV
    \placetable[][tab:#1]{#4}{%
      \processseparatedfile[MarkdownConTeXtCSV][#3]}%
  \else
    \markdownInput{#3}%
  \fi}%
\def\markdownRendererImagePrototype#1#2#3#4{%
  \placefigure[][]{#4}{\externalfigure[#3]}}%
\def\markdownRendererUlBeginPrototype{\startitemize}%
\def\markdownRendererUlBeginTightPrototype{\startitemize[packed]}%
\def\markdownRendererUlItemPrototype{\item}%
\def\markdownRendererUlEndPrototype{\stopitemize}%
\def\markdownRendererUlEndTightPrototype{\stopitemize}%
\def\markdownRendererOlBeginPrototype{\startitemize[n]}%
\def\markdownRendererOlBeginTightPrototype{\startitemize[packed,n]}%
\def\markdownRendererOlItemPrototype{\item}%
\def\markdownRendererOlItemWithNumberPrototype#1{\sym{#1.}}%
\def\markdownRendererOlEndPrototype{\stopitemize}%
\def\markdownRendererOlEndTightPrototype{\stopitemize}%
\definedescription
  [MarkdownConTeXtDlItemPrototype]
  [location=hanging,
   margin=standard,
   headstyle=bold]%
\definestartstop
  [MarkdownConTeXtDlPrototype]
  [before=\blank,
   after=\blank]%
\definestartstop
  [MarkdownConTeXtDlTightPrototype]
  [before=\blank\startpacked,
   after=\stoppacked\blank]%
\def\markdownRendererDlBeginPrototype{%
  \startMarkdownConTeXtDlPrototype}%
\def\markdownRendererDlBeginTightPrototype{%
  \startMarkdownConTeXtDlTightPrototype}%
\def\markdownRendererDlItemPrototype#1{%
  \startMarkdownConTeXtDlItemPrototype{#1}}%
\def\markdownRendererDlItemEndPrototype{%
  \stopMarkdownConTeXtDlItemPrototype}%
\def\markdownRendererDlEndPrototype{%
  \stopMarkdownConTeXtDlPrototype}%
\def\markdownRendererDlEndTightPrototype{%
  \stopMarkdownConTeXtDlTightPrototype}%
\def\markdownRendererEmphasisPrototype#1{{\em#1}}%
\def\markdownRendererStrongEmphasisPrototype#1{{\bf#1}}%
\def\markdownRendererBlockQuoteBeginPrototype{\startquotation}%
\def\markdownRendererBlockQuoteEndPrototype{\stopquotation}%
\def\markdownRendererInputVerbatimPrototype#1{\typefile{#1}}%
\def\markdownRendererInputFencedCodePrototype#1#2{%
  \ifx\relax#2\relax
    \typefile{#1}%
  \else
    \typefile[#2][]{#1}%
  \fi}%
\def\markdownRendererHeadingOnePrototype#1{\chapter{#1}}%
\def\markdownRendererHeadingTwoPrototype#1{\section{#1}}%
\def\markdownRendererHeadingThreePrototype#1{\subsection{#1}}%
\def\markdownRendererHeadingFourPrototype#1{\subsubsection{#1}}%
\def\markdownRendererHeadingFivePrototype#1{\subsubsubsection{#1}}%
\def\markdownRendererHeadingSixPrototype#1{\subsubsubsubsection{#1}}%
\def\markdownRendererHorizontalRulePrototype{%
  \blackrule[height=1pt, width=\hsize]}%
\def\markdownRendererFootnotePrototype#1{\footnote{#1}}%
\def\markdownRendererTickedBoxPrototype{$\boxtimes$}
\def\markdownRendererHalfTickedBoxPrototype{$\boxdot$}
\def\markdownRendererUntickedBoxPrototype{$\square$}
\def\markdownRendererStrikeThroughPrototype#1{\overstrikes{#1}}
\def\markdownRendererSuperscriptPrototype#1{\high{#1}}
\def\markdownRendererSubscriptPrototype#1{\low{#1}}
\newcount\markdownConTeXtRowCounter
\newcount\markdownConTeXtRowTotal
\newcount\markdownConTeXtColumnCounter
\newcount\markdownConTeXtColumnTotal
\newtoks\markdownConTeXtTable
\newtoks\markdownConTeXtTableFloat
\def\markdownRendererTablePrototype#1#2#3{%
  \markdownConTeXtTable={}%
  \ifx\empty#1\empty
    \markdownConTeXtTableFloat={%
      \the\markdownConTeXtTable}%
  \else
    \markdownConTeXtTableFloat={%
      \placetable{#1}{\the\markdownConTeXtTable}}%
  \fi
  \begingroup
  \setupTABLE[r][each][topframe=off, bottomframe=off, leftframe=off, rightframe=off]
  \setupTABLE[c][each][topframe=off, bottomframe=off, leftframe=off, rightframe=off]
  \setupTABLE[r][1][topframe=on, bottomframe=on]
  \setupTABLE[r][#1][bottomframe=on]
  \markdownConTeXtRowCounter=0%
  \markdownConTeXtRowTotal=#2%
  \markdownConTeXtColumnTotal=#3%
  \markdownConTeXtRenderTableRow}
\def\markdownConTeXtRenderTableRow#1{%
  \markdownConTeXtColumnCounter=0%
  \ifnum\markdownConTeXtRowCounter=0\relax
    \markdownConTeXtReadAlignments#1%
    \markdownConTeXtTable={\bTABLE}%
  \else
    \markdownConTeXtTable=\expandafter{%
      \the\markdownConTeXtTable\bTR}%
    \markdownConTeXtRenderTableCell#1%
    \markdownConTeXtTable=\expandafter{%
      \the\markdownConTeXtTable\eTR}%
  \fi
  \advance\markdownConTeXtRowCounter by 1\relax
  \ifnum\markdownConTeXtRowCounter>\markdownConTeXtRowTotal\relax
    \markdownConTeXtTable=\expandafter{%
      \the\markdownConTeXtTable\eTABLE}%
    \the\markdownConTeXtTableFloat
    \endgroup
    \expandafter\gobbleoneargument
  \fi\markdownConTeXtRenderTableRow}
\def\markdownConTeXtReadAlignments#1{%
  \advance\markdownConTeXtColumnCounter by 1\relax
  \if#1d%
    \setupTABLE[c][\the\markdownConTeXtColumnCounter][align=right]
  \fi\if#1l%
    \setupTABLE[c][\the\markdownConTeXtColumnCounter][align=right]
  \fi\if#1c%
    \setupTABLE[c][\the\markdownConTeXtColumnCounter][align=middle]
  \fi\if#1r%
    \setupTABLE[c][\the\markdownConTeXtColumnCounter][align=left]
  \fi
  \ifnum\markdownConTeXtColumnCounter<\markdownConTeXtColumnTotal\relax\else
    \expandafter\gobbleoneargument
  \fi\markdownConTeXtReadAlignments}
\def\markdownConTeXtRenderTableCell#1{%
  \advance\markdownConTeXtColumnCounter by 1\relax
  \markdownConTeXtTable=\expandafter{%
    \the\markdownConTeXtTable\bTD#1\eTD}%
  \ifnum\markdownConTeXtColumnCounter<\markdownConTeXtColumnTotal\relax\else
    \expandafter\gobbleoneargument
  \fi\markdownConTeXtRenderTableCell}
\stopmodule\protect
\endinput
%%
%% End of file `t-markdown.tex'.
